% Generated draft: requires mathematical and source review.
\documentclass[11pt]{article}
\usepackage[margin=25mm]{geometry}
\usepackage{fontspec}
\setmainfont{Latin Modern Roman}
\usepackage{amsmath,amssymb,mathtools,unicode-math}
\setmathfont{Latin Modern Math}
\usepackage{enumitem,needspace,xurl,hyperref}
\setlength{\parindent}{0pt}
\setlength{\parskip}{5pt}
\setlength{\emergencystretch}{3em}
\newcommand{\N}{\mathbb{N}}
\newcommand{\Z}{\mathbb{Z}}
\newcommand{\Q}{\mathbb{Q}}
\newcommand{\R}{\mathbb{R}}
\newcommand{\C}{\mathbb{C}}
\newcommand{\F}{\mathbb{F}}
\newcommand{\A}{\mathbb{A}}
\newcommand{\PP}{\mathbb{P}}
\newcommand{\E}{\mathbb{E}}
\newcommand{\eps}{\varepsilon}
\newcommand{\dd}{\,\mathrm{d}}
\begin{document}
\section*{UnsolvedMath ID: 30006165}
\textbf{Source problem code:} catalogue entry
\subsection{Definitions and mathematical statement}
A standard Borel space is a measurable space isomorphic to the Borel measurable space of a Polish space. For equivalence relations \(E\) on a standard Borel space \(X\) and \(F\) on a standard Borel space \(Y\), a Borel reduction of \(E\) to \(F\) is a Borel map \(f:X\to Y\) satisfying
\[
\forall x,x'\in X\qquad x\,E\,x'\iff f(x)\,F\,f(x').
\]
Write \(E\leq_B F\) when such a map exists, and \(E\sim_B F\) when both \(E\leq_B F\) and \(F\leq_B E\).

Under the usual descriptive set theoretic convention, an orbit equivalence relation is obtained from a Polish group \(G\) acting on a standard Borel space \(Y\) through a Borel action map \(a:G\times Y\to Y\). Here \(a(1_G,y)=y\) and \(a(gh,y)=a(g,a(h,y))\). The associated relation is
\[
y\,E_a\,y'\iff (\exists g\in G)\;a(g,y)=y'.
\]
The supplied record does not specify this acting-group and action convention, so its agreement with the source must be checked.

The supplied record also does not define \emph{idealistic}. Let \(\operatorname{Ideal}_{\mathrm{src}}(E)\) stand for the intended source property; this is an unresolved predicate, not a replacement definition. In particular, no additional assumption that \(E\) is a Borel subset of \(X\times X\) is inserted.

Subject to resolving that predicate and confirming the orbit convention, the question is whether
\[
\begin{gathered}
\text{for every standard Borel space }X\text{ and every equivalence relation }E\text{ on }X,\\
\operatorname{Ideal}_{\mathrm{src}}(E)\ \Longrightarrow\
\exists\,Y,G,a,f,h\quad
\left\{
\begin{aligned}
&Y\text{ is a standard Borel space},\quad G\text{ is a Polish group},\\
&a:G\times Y\to Y\text{ is a Borel action},\\
&f:X\to Y\text{ and }h:Y\to X\text{ are Borel},\\
&\forall x,x'\in X\quad
x\,E\,x'\iff f(x)\,E_a\,f(x'),\\
&\forall y,y'\in Y\quad
y\,E_a\,y'\iff h(y)\,E\,h(y').
\end{aligned}
\right.
\end{gathered}
\]
Thus the requested conclusion is \(E\sim_B E_a\). The missing definition of idealistic prevents this draft from being a fully self-contained formulation. The supplied source labels the question open; that status is not independently verified here.
\subsection{Short English statement}
The question asks whether every idealistic equivalence relation on a standard Borel space has exactly the same complexity under Borel reducibility as some orbit equivalence relation. The supplied record leaves the definition of idealistic and the precise orbit-action convention unstated; these require source confirmation.
\subsection{Sources}
\begin{itemize}
\item Supplied local UnsolvedMath registry record, permanent ID \(30006165\), title \emph{Idealistic Equivalence Relations and Orbit Relations}, exact source problem code \texttt{catalogue entry}, in the Oberwolfach Reports Open Problems collection.
\item \emph{Set Theory} (2025), Oberwolfach Reports, DOI \texttt{10.4171/owr/2025/2}, pp.\ 75--140. Supplied official-report reference: \url{https://doi.org/10.4171/owr/2025/2}. The excerpt attributes \emph{Idealistic equivalence relations remastered} to Filippo Calderoni, joint work with Luca Motto Ros. The record gives 2025 as the proposed year but leaves its separate proposer field empty.
\item Calderoni--Motto Ros, \emph{Structural results on idealistic equivalence relations}, supplied manuscript link: \url{https://sites.math.rutgers.edu/~fc327/idealistic.pdf}.
\item Literature assessment dated exactly \texttt{2026-08-22}: the supplied curation labels the question open and reports that the 2026 primary manuscript still describes the general Borel bireducibility question as open, while proving a strong obstruction for a related classwise embedding formulation. This is supplied dated curation, not an independently verified assessment; its time-sensitive status and cited claims require independent checking.
\item The original excerpt mentions the reference label \texttt{[KecLou]}, without supplying its bibliographic entry. Its identity and relevance remain unresolved.
\item The registry labels the statement \texttt{corrected\_verified} and says it was reconstructed and checked against the official Oberwolfach report. This records the supplied verification claim; no source was independently visited or checked for this draft.
\end{itemize}
\end{document}
